% Generated draft: requires mathematical and source review.
\documentclass[11pt]{article}
\usepackage[margin=25mm]{geometry}
\usepackage{fontspec}
\setmainfont{Latin Modern Roman}
\usepackage{amsmath,amssymb,mathtools,unicode-math}
\setmathfont{Latin Modern Math}
\usepackage{enumitem,needspace,xurl,hyperref}
\setlength{\parindent}{0pt}
\setlength{\parskip}{5pt}
\setlength{\emergencystretch}{3em}
\newcommand{\N}{\mathbb{N}}
\newcommand{\Z}{\mathbb{Z}}
\newcommand{\Q}{\mathbb{Q}}
\newcommand{\R}{\mathbb{R}}
\newcommand{\C}{\mathbb{C}}
\newcommand{\F}{\mathbb{F}}
\newcommand{\A}{\mathbb{A}}
\newcommand{\PP}{\mathbb{P}}
\newcommand{\E}{\mathbb{E}}
\newcommand{\eps}{\varepsilon}
\newcommand{\dd}{\,\mathrm{d}}
\begin{document}
\section*{UnsolvedMath ID: 800853}
\textbf{Source problem code:} catalogue entry
\subsection{Definitions and mathematical statement}
For a prime \(p\equiv 1\pmod 4\), let \(\F_p=\Z/p\Z\). The Paley graph \(G_p\) is the simple undirected graph with vertex set \(\F_p\), in which distinct vertices \(x,y\) are adjacent if \(x-y\) is a nonzero square in \(\F_p\). The congruence condition ensures that adjacency is symmetric. Define its clique number by
\[
\omega(G_p)=\max\bigl\{|S|:S\subseteq\F_p,\ x-y\text{ is a nonzero square for all distinct }x,y\in S\bigr\}.
\]
The independence number \(\alpha(G_p)\) is the largest size of a vertex set containing no adjacent pair. As noted in the supplied source, \(G_p\) is isomorphic to its complement, so \(\omega(G_p)=\alpha(G_p)\).

The supplied Conjecture 1 asks whether
\[
\omega(G_p)=O\bigl((\log p)^k\bigr)
\]
for some fixed exponent \(k>0\), as \(p\) tends to infinity through primes congruent to \(1\) modulo \(4\). Explicitly, the conjecture is that there exist constants \(C>0\), \(k>0\), and \(p_0\geq 2\) such that for every prime \(p\geq p_0\) with \(p\equiv1\pmod4\),
\[
\omega(G_p)\leq C(\log p)^k.
\]
Here \(\log\) denotes the natural logarithm; the constants are independent of \(p\). This interprets the source's unspecified \(\mathrm{polylog}(p)\) as a fixed power of \(\log p\); no particular exponent is supplied. The source motivates this conjecture by comparison with a random graph of the same degree, but does not specify the random graph model. That comparison is background motivation and imposes no additional hypothesis on the conjecture.
\subsection{Short English statement}
The conjecture asks whether the largest clique in the Paley graph on \(p\) vertices is bounded by a fixed power of \(\log p\), uniformly over primes \(p\equiv1\pmod4\). By self-complementarity, the same bound would hold for the independence number. The registry labels the problem open; this is supplied status information, not an independently verified literature assessment.
\subsection{Sources}
\begin{itemize}
\item Local UnsolvedMath registry, supplied complete record: permanent ID \(800853\), source problem code \texttt{catalogue entry}, title ``Randomstrasse101 Open Problems (2025) --- Conjecture 1.'' The registry attributes the record to Bandeira et al., gives proposed year 2026, and labels it open.
\item Bandeira et al., \emph{Randomstrasse101: Open Problems of 2025} (2026), Conjecture 1. Supplied source URL: \url{https://arxiv.org/abs/2603.29571}. The supplied provenance records access on 2026-10-16 and extraction by \texttt{arxiv\_latex}; these are supplied provenance claims, not a claim of independent access or verification.
\item The extracted statement describes the conjecture as well known and points to ``Open Problem 8.4 in [Afonso10L42P].'' The supplied record does not resolve that citation key into a bibliographic reference or URL, so the earlier attribution cannot be specified further from the supplied evidence.
\end{itemize}
\end{document}
